\section*{Discussion}

This tissue-aware screen produced two main findings. First, the threshold-defined primary analysis recovered TYK2 as the only FDR-significant gene, with psoriasis associations in both skin tissues and whole blood. Second, strict instrument coverage was limited: only 18 of 57 gene--tissue cells had a genome-wide-significant cis-eQTL, and only 17 cells could be matched to the four FinnGen outcomes. The study therefore functions more convincingly as an assessment of genetic tractability within the curated panel and an established-locus positive control than as a discovery screen for new drug targets.

The TYK2 result requires careful directional interpretation. The eQTL alleles associated with higher measured TYK2 transcript were associated with lower psoriasis risk. This should not be translated into a claim that increasing TYK2 expression is therapeutic or that the eQTL directly mimics pharmacologic inhibition. The TYK2 locus contains coding and regulatory variation, transcript abundance may not proxy kinase activity, and bulk-tissue eQTL effects may differ across cell types or activation states. The MR odds ratios are scaled to genetically proxied transcript differences and are not estimates of deucravacitinib or any other drug effect. The value of TYK2 as a positive control is therefore the recovery and validation of an established psoriasis locus, not a simple match between the sign of the transcript estimate and the direction of deucravacitinib treatment.\textsuperscript{\cite{steri2025tyk2,armstrong2023deucravacitinib}}

The colocalization results refine that interpretation. Strong shared-variant support in sun-exposed skin and supportive evidence in whole blood strengthen the case that the expression and psoriasis signals overlap in those datasets. In non-sun-exposed skin, PP.H2 exceeded PP.H4, so the MR association alone is insufficient to infer shared causality. Colocalization itself does not prove that altered TYK2 expression is the mediating molecular mechanism, particularly in a locus with coding variation, but it reduces concern that the two associations arise from completely distinct variants in sun-exposed skin and whole blood.\textsuperscript{\cite{giambartolomei2014coloc,zuber2022coloc}}

Instrument scarcity was not confined to one pathway. At the primary threshold, only nine of 19 genes were represented in at least one tissue. A relaxed threshold added IL36RN and PDE4B, raising gene-level coverage to 11 of 19, but eight genes remained uninstrumented. Several are clinically important, including TNF, IL17A, IL23R, and JAK3. Their absence from the MR results reflects the limitations of bulk-tissue expression resources, the selected thresholds, and the use of GTEx significant-pair rather than complete all-pairs files; it is not evidence against their established therapeutic relevance.\textsuperscript{\cite{gill2024pitfalls}} Future work may improve coverage with larger skin eQTL studies, cell-type-specific QTLs, coding variants, or rigorously curated protein QTL resources.\textsuperscript{\cite{sun2023plasma,ferkingstad2021proteome,trajanoska2023target}}

The separate PDE4A result illustrates why threshold provenance and colocalization matter. Its association was statistically strong within the 20-test relaxed family, but the exposure eQTL did not meet the primary genome-wide-significance criterion and locus-level PP.H4 was only 0.108. The selected MR instrument was absent from the eQTL Catalogue rows used for colocalization, so the analysis did not directly test that variant and is non-supportive rather than definitive. The result may reflect instrument-selection uncertainty, LD, or locus architecture and should not be presented as genetic validation of PDE4A inhibition. JAK2 was not testable because no eligible row met either threshold; accordingly, no JAK2 estimate is reported.

\subsection*{Clinical Implications for Target Selection}

The findings define three clinically distinct evidence states. First, recovery of an established locus with shared-variant support, as for TYK2 in sun-exposed skin and whole blood, justifies further mechanistic study but does not by itself establish the direction or magnitude of a pharmacologic intervention. Second, absence of an eligible instrument, as for TNF, IL17A, IL23R, and JAK3, means that the target is not testable in this dataset and should be judged using other evidence streams, including disease-tissue biology, functional studies, and clinical trials. Third, an association based on a relaxed instrument without colocalization support, as for PDE4A, should remain a low-confidence hypothesis until replicated with stronger instruments, fine-mapping, or orthogonal molecular data. This three-state interpretation is more useful for clinical target selection than a binary supported-versus-unsupported label.

This study has several strengths. All seven compressed source files were reread, the 68 primary estimates were independently reproduced to numerical precision, primary and relaxed hypotheses were corrected separately, and every reported MR row can be traced to its exposure and outcome beta, standard error, and allele orientation. Reporting the 57-cell availability matrix also makes the denominator of the screen explicit. Finally, tissue-resolved colocalization prevented a uniform interpretation of three superficially concordant TYK2 MR estimates.

The limitations are substantial. Each estimate used one variant, so MR-Egger, weighted median, heterogeneity testing, and MR-PRESSO were not possible. No external LD reference was available, and lead-variant selection is not a substitute for LD clumping or correlated-instrument cis-MR.\textsuperscript{\cite{gkatzionis2022cismr,lin2024cismr}} A single-SNP Wald ratio cannot exclude horizontal pleiotropy, whose direction is not predictable from the available data. The three TYK2 tissue rows rely on only two distinct variants and are partly non-independent. The significant-pair GTEx files do not provide a complete inventory of nominal cis associations, so the relaxed-threshold coverage is conditional and may be underestimated. The single-causal-variant assumption of \textit{coloc.abf} may be violated by allelic heterogeneity, which could affect the reported posterior probabilities. The eQTL colocalization inputs also used $sdY=1$, an assumed unit phenotypic standard deviation whose misspecification could affect Bayes factors and posterior probabilities. Reverse MR and Steiger analyses could not be validly performed with significant-pair-only GTEx files and unverified tissue sample sizes. Formal power percentages were not defensible; the much smaller vitiligo and alopecia areata case counts imply lower precision and warrant caution when interpreting null results. The study was not prospectively registered. GTEx represents mostly healthy bulk tissue, FinnGen is ancestry-specific, and the external GWAS records may overlap. Lifelong genetically proxied expression differences cannot be assumed to represent treatment timing, dose, duration, or expression changes outside the genetically observed range. The target panel also mixes direct targets with pathway and disease genes, so results should be interpreted at the gene and locus level rather than as uniform pharmacologic proxies.

Within these constraints, the data support a narrow conclusion: the pipeline recovers the established TYK2 psoriasis locus, with the strongest shared-variant evidence in sun-exposed skin and whole blood, while current GTEx resources leave most panel gene--tissue cells untestable. Dermatologic drug-target MR studies should therefore report instrument coverage, threshold provenance, and colocalization alongside effect estimates.
